\begin{lemma}
Let $n\ge0$ be an integer, let $0\le a<b$, let $k\in\mathbb{Z}$, put
\[
  A=\{\xi\in\mathbb{C}:a<|\xi|<b\},
\]
and let $H,D_0,D_1\colon\mathbb{C}\to\mathbb{C}$ satisfy
\begin{enumerate}
\item[(H1)] $H\in L^{1}_{\mathrm{loc}}(A)$ and $D_j\in L^{1}_{\mathrm{loc}}(A)$
  for $j=0,1$, in the sense of \noderef{MemLpLocOn};
\item[(H2)] $D_j$ is the distributional $j$-th partial derivative of $H$ on
  $A$ for $j=0,1$, in the sense of \noderef{IsWeakPartialOn};
\item[(H3)] $H$ is continuous on $A$;
\item[(H4)] $H$ is differentiable at almost every point of $A$.
\end{enumerate}
Write
\[
  a_k(r)=\frac1{2\pi}\int_0^{2\pi}H(re^{i\theta})e^{i(n-k)\theta}\,d\theta,
  \qquad
  A_k(r)=\frac1{2\pi}\int_0^{2\pi}H_r(re^{i\theta})e^{i(n-k)\theta}\,d\theta,
  \qquad
  B_k(r)=\frac1{2\pi}\int_0^{2\pi}H_\theta(re^{i\theta})e^{i(n-k)\theta}\,d\theta
\]
where the three displayed expressions are the Fourier coefficient (30) of
\noderef{FourierCoefOfH} formed from $H$, from the polar radial derivative
$H_r$ of \noderef{HRad} and from the polar angular derivative $H_\theta$ of
\noderef{HAng} respectively.  Then
\begin{enumerate}
\item[(a)] for all $c,d\in(a,b)$, $A_k$ is integrable on the closed interval
  with endpoints $c$ and $d$;
\item[(b)] for all $c,d\in(a,b)$,
  $a_k(d)-a_k(c)=\displaystyle\int_c^dA_k(s)\,ds$;
\item[(c)] for almost every $r\in(a,b)$, $a_k$ is differentiable at $r$ with
  $a_k'(r)=A_k(r)$;
\item[(d)] for almost every $r\in(a,b)$, $B_k(r)=i(k-n)\,a_k(r)$;
\item[(e)] for almost every $r\in(a,b)$, the functions
  $\theta\mapsto H_r(re^{i\theta})$ and $\theta\mapsto H_\theta(re^{i\theta})$
  are integrable on $[0,2\pi]$.
\end{enumerate}

This is the one-variable content of the paper's (31) and of the two identities
that the term-by-term differentiation (32) amounts to.  It is what replaces
Section 3's standing assumption $H\in C^\infty(\mathbb{B}-\{0\})$: the paper
reads that assumption at exactly one place, paper.tex:265--266, to obtain
$a_k\in C[0,1)\cap C^\infty(0,1)$, while paper.tex:269 justifies (32) from
$H\in W^{1,2}(\mathbb{B})$ alone; and the $C[0,1)$ half of (31) is continuity,
which the inverse of a quasiconformal homeomorphism has for nothing.  (a), (b)
and (c) say that $a_k$ is locally absolutely continuous on $(a,b)$ in the
strongest usable sense -- it is the indefinite integral of $A_k$ -- and (d) is
the integration by parts in $\theta$ that the paper performs at
paper.tex:281--286 to turn (28) into the system (33).  The hypotheses are
exactly the data of a $W^{1,1}_{\mathrm{loc}}(A)$ function, with $D_j$ the
member of the coherent family indexed by the one-letter word $[j]$, together
with the continuity and almost-everywhere differentiability that item (C2) of
the paper's note supplies for the inverse of a quasiconformal homeomorphism.
That the paper means its circle averages to be differentiated in exactly this
way is explicit at paper.tex:268--277, where $H\in W^{1,2}(\mathbb{B})$ alone
is what justifies (32) and its two series are said to converge in
$L^{2}(\mathbb{B})$, the derivatives $H_r$ and $H_\theta$ appearing there being
the classical ones.

The conclusions are no stronger than that: the derivative identity can only
hold almost everywhere, since $H_r$ and $H_\theta$ are defined through the
Fr\'echet derivative of $H$, which (H4) provides only almost everywhere.  They
are also deliberately short of (31) itself.  Nothing here concerns the centre
of the disk, or the family of all $k$ at once, or the $L^{2}$ expansion (29),
or $C^\infty$ regularity in $r$: the $C^\infty$ half of (31) is a consequence
of the system (33) satisfied by the coefficients, not of this lemma, and the
$C[0,1)$ half is continuity.
\end{lemma}
\begin{proof}
Write $d\sigma$ for plane Lebesgue measure, $e_0=1$, $e_1=i$, and
$\partial_jv(\zeta)=\mathrm{d}v(\zeta)[e_j]$.  The set $A$ is open, and
$0\notin A$ because $|\xi|>a\ge0$ for $\xi\in A$.  For $\zeta\in\mathbb{C}$
write $\zeta=\zeta_1+i\zeta_2$ with $\zeta_1,\zeta_2$ real.

\emph{Step 1: the radial and angular combinations of the weak partials.}
Define, for $\zeta\neq0$,
\[
  R(\zeta)=\frac{\zeta_1}{|\zeta|}D_0(\zeta)+\frac{\zeta_2}{|\zeta|}D_1(\zeta),
  \qquad
  S(\zeta)=-\zeta_2D_0(\zeta)+\zeta_1D_1(\zeta).
\]
The hypotheses of \noderef{WeakPartialEqFDerivAE} hold with $\Omega=A$, which
is open, the function $H$, the family $D$, and (H1), (H2), (H4).  Its
conclusion gives, for each of $j=0$ and $j=1$, a null set outside which
$D_j=\partial_jH$ on $A$; let $N_1$ be the union of those two null sets and of
the null set outside which (H4) makes $H$ differentiable on $A$, so that $N_1$
is null and, for every $\zeta\in A-N_1$, $H$ is differentiable at $\zeta$ and
$D_j(\zeta)=\partial_jH(\zeta)$ for $j=0,1$.  Fix
$\zeta\in A-N_1$ and put $r=|\zeta|>0$.  Since
$\zeta/|\zeta|=(\zeta_1/r)e_0+(\zeta_2/r)e_1$ and
$i\zeta=-\zeta_2e_0+\zeta_1e_1$, and $\mathrm{d}H(\zeta)$ is
$\mathbb{R}$-linear, \noderef{HRad} and \noderef{HAng} give
\begin{equation}\tag{1}
  H_r(\zeta)=\mathrm{d}H(\zeta)\Bigl[\frac{\zeta}{|\zeta|}\Bigr]=R(\zeta),
  \qquad
  H_\theta(\zeta)=\mathrm{d}H(\zeta)[i\zeta]=S(\zeta)
  \qquad(\zeta\in A-N_1).
\end{equation}

\emph{Step 2: polar consequences of (1).}  Let $[c',d']\subseteq(a,b)$ be a
closed interval and put
$A'=\{\zeta:c'\le|\zeta|\le d'\}$, a compact subset of $A$.  On $A'$ the
coefficients of $R$ and $S$ are bounded, by $1$ and by $d'$ respectively, so by
(H1) applied to the compact set $A'$ both $R$ and $S$ are integrable on $A'$.
Polar coordinates, that is the change of variables $(s,\theta)\mapsto
se^{i\theta}$ under which $d\sigma=s\,ds\,d\theta$, together with Tonelli's
theorem and $s\ge c'>0$ on $A'$, give
\begin{equation}\tag{2}
  \int_{c'}^{d'}\!\!\int_0^{2\pi}|R(se^{i\theta})|\,d\theta\,ds
  \le\frac1{c'}\int_{A'}|R|\,d\sigma<\infty,
\end{equation}
and the same bound for $S$.  This estimate, together with the two consequences
drawn from it below --- that the coefficient of a function integrable on the
annulus is integrable in the radius, and that the coefficient depends
continuously on the function in $L^1$ --- is
\noderef{FourierCoefAnnulusRegularity}, stated there for an arbitrary
measurable function on a compact annulus of positive inner radius, and it is
what Step 4's mollification limit is applied through.  By Fubini's theorem the
function
$s\mapsto\frac1{2\pi}\int_0^{2\pi}R(se^{i\theta})e^{i(n-k)\theta}\,d\theta$ is
therefore defined for almost every $s\in[c',d']$, measurable, and integrable
there; write $\widetilde A_k$ for it, and $\widetilde B_k$ for the same
expression with $S$ in place of $R$.  Moreover the null set $N_1$ of Step 1 has
$\int_{c'}^{d'}\int_0^{2\pi}\mathbf1_{N_1}(se^{i\theta})\,s\,d\theta\,ds=0$,
so for almost every $s\in[c',d']$ the circle $\{se^{i\theta}:0\le\theta<2\pi\}$
meets $N_1$ in a set of $\theta$-measure zero; for such $s$ the integrands of
$A_k(s)$ and $\widetilde A_k(s)$ agree for almost every $\theta$ by (1), as do
those of $B_k(s)$ and $\widetilde B_k(s)$.  For almost every
$s\in[c',d']$ the integrand of $\widetilde A_k(s)$ is moreover integrable in
$\theta$, by the Fubini theorem just invoked, hence so is the integrand of
$A_k(s)$, two functions agreeing almost everywhere being integrable together;
and then the two integrals coincide.  The same applies to $B_k$ and
$\widetilde B_k$.  Hence
\begin{equation}\tag{3}
  A_k(s)=\widetilde A_k(s)\quad\text{and}\quad B_k(s)=\widetilde B_k(s)
  \qquad\text{for almost every }s\in[c',d'] .
\end{equation}
In particular $A_k$ agrees almost everywhere on $[c',d']$ with an integrable
function, so $A_k$ is integrable on $[c',d']$.  Since every $c,d\in(a,b)$ lie
in such an interval $[c',d']\subseteq(a,b)$, this proves (a).

\emph{Step 3: the mollification.}  Fix $c,d\in(a,b)$ and assume first that
$c\le d$; the opposite case is reduced to this one at the end of Step 6.  Choose $\delta>0$ with
$[c-3\delta,d+3\delta]\subseteq(a,b)$, which is possible because $(a,b)$ is
open and $[c,d]\subseteq(a,b)$ is compact; note $c-3\delta>a\ge0$.  Put
\[
  A''=\{\zeta:c-3\delta\le|\zeta|\le d+3\delta\},\qquad
  A^{\flat}=\{\zeta:c-\delta\le|\zeta|\le d+\delta\},
\]
both compact subsets of $A$, with $A^\flat\subseteq A''$.  Fix
$\varrho\colon\mathbb{C}\to\mathbb{R}$ of class $C^\infty$ with $\varrho\ge0$,
$\varrho(y)=0$ for $|y|\ge1$ and $\int\varrho\,d\sigma=1$, and for
$0<\varepsilon\le\delta$ put
$\varrho_\varepsilon(y)=\varepsilon^{-2}\varrho(y/\varepsilon)$, so that
$\varrho_\varepsilon$ is $C^\infty$, vanishes for $|y|\ge\varepsilon$ and has
integral $1$.  By (H1) applied to the compact set $A''$ the functions
$\mathbf1_{A''}H$ and $\mathbf1_{A''}D_j$ lie in $L^{1}(\mathbb{C})$; put
\[
  H_\varepsilon=(\mathbf1_{A''}H)*\varrho_\varepsilon,\qquad
  W_{j,\varepsilon}=(\mathbf1_{A''}D_j)*\varrho_\varepsilon .
\]
These are $C^\infty$ on $\mathbb{C}$, being convolutions of $L^{1}$ functions
with a compactly supported $C^\infty$ function, and
$\partial_jH_\varepsilon=(\mathbf1_{A''}H)*\partial_j\varrho_\varepsilon$.

\emph{Step 4: $\partial_jH_\varepsilon=W_{j,\varepsilon}$ on $A^\flat$.}  Let
$\zeta\in A^\flat$ and $0<\varepsilon\le\delta$, and put
$\varphi(w)=\varrho_\varepsilon(\zeta-w)$, read as a complex-valued function.
Then $\varphi$ is $C^\infty$ and vanishes off
$\{w:|w-\zeta|\le\varepsilon\}$, a compact set contained in $A''\subseteq A$,
because $|\,|w|-|\zeta|\,|\le\varepsilon\le\delta$ there and
$c-\delta\le|\zeta|\le d+\delta$.  So
$\operatorname{tsupp}\varphi\subseteq A$ and $\varphi$ is a test function for
$A$ in the sense of \noderef{IsTestFn}.  The map $w\mapsto\zeta-w$ is affine
with derivative $-\mathrm{id}$, so
$\mathrm{d}\varphi(w)[e_j]=-(\partial_j\varrho_\varepsilon)(\zeta-w)$, and (H2)
applied to $\varphi$ reads
\[
  -\int H(w)(\partial_j\varrho_\varepsilon)(\zeta-w)\,d\sigma(w)
  =-\int D_j(w)\varrho_\varepsilon(\zeta-w)\,d\sigma(w).
\]
Both integrands vanish off $\{w:|w-\zeta|\le\varepsilon\}\subseteq A''$, where
$H=\mathbf1_{A''}H$ and $D_j=\mathbf1_{A''}D_j$, so the displayed identity says
precisely
\begin{equation}\tag{4}
  \partial_jH_\varepsilon(\zeta)=W_{j,\varepsilon}(\zeta)
  \qquad(\zeta\in A^\flat,\ 0<\varepsilon\le\delta).
\end{equation}
Consequently, with
$R_\varepsilon(\zeta)=\frac{\zeta_1}{|\zeta|}W_{0,\varepsilon}(\zeta)
+\frac{\zeta_2}{|\zeta|}W_{1,\varepsilon}(\zeta)$ and
$S_\varepsilon(\zeta)=-\zeta_2W_{0,\varepsilon}(\zeta)
+\zeta_1W_{1,\varepsilon}(\zeta)$, the computation of Step 1 applied to the
everywhere differentiable $H_\varepsilon$ gives
\begin{equation}\tag{5}
  (H_\varepsilon)_r=R_\varepsilon\quad\text{and}\quad
  (H_\varepsilon)_\theta=S_\varepsilon\qquad\text{on }A^\flat .
\end{equation}

\emph{Step 5: the radial identity for $H_\varepsilon$.}  Put
\[
  g_\varepsilon(r)=\frac1{2\pi}\int_0^{2\pi}
    H_\varepsilon(re^{i\theta})e^{i(n-k)\theta}\,d\theta .
\]
The function $(r,\theta)\mapsto H_\varepsilon(re^{i\theta})e^{i(n-k)\theta}$ is
$C^\infty$ on $\mathbb{R}^2$ and the $\theta$-range is a compact interval, so
differentiation under the integral sign applies: $g_\varepsilon$ is $C^1$ with
\[
  g_\varepsilon'(r)=\frac1{2\pi}\int_0^{2\pi}
    \frac{\partial}{\partial r}\bigl[H_\varepsilon(re^{i\theta})\bigr]
    e^{i(n-k)\theta}\,d\theta .
\]
The curve $r\mapsto re^{i\theta}$ has velocity $e^{i\theta}$, so by the chain
rule $\frac{\partial}{\partial r}[H_\varepsilon(re^{i\theta})]
=\mathrm{d}H_\varepsilon(re^{i\theta})[e^{i\theta}]$, and for $r>0$ the vector
$e^{i\theta}$ is $\xi/|\xi|$ at $\xi=re^{i\theta}$, so this is
$(H_\varepsilon)_r(re^{i\theta})$ by \noderef{HRad}.  Hence, for $r>0$,
$g_\varepsilon'(r)=\frac1{2\pi}\int_0^{2\pi}(H_\varepsilon)_r(re^{i\theta})
e^{i(n-k)\theta}\,d\theta$, and the fundamental theorem of calculus for the
$C^1$ function $g_\varepsilon$ on $[c,d]\subseteq(0,\infty)$ gives, using (5)
and $\{se^{i\theta}:c\le s\le d\}\subseteq A^\flat$,
\begin{equation}\tag{6}
  g_\varepsilon(d)-g_\varepsilon(c)
  =\int_c^d\Bigl(\frac1{2\pi}\int_0^{2\pi}
     R_\varepsilon(se^{i\theta})e^{i(n-k)\theta}\,d\theta\Bigr)ds .
\end{equation}

\emph{Step 6: letting $\varepsilon\to0$ in (6); proof of (b).}  First,
$H_\varepsilon\to H$ uniformly on $A^\flat$.  Indeed $H$ is continuous on $A$
by (H3) and $A''\subseteq A$ is compact, so $H$ is uniformly continuous on
$A''$; and for $\zeta\in A^\flat$,
\[
  H_\varepsilon(\zeta)-H(\zeta)
  =\int\bigl[H(\zeta-y)-H(\zeta)\bigr]\varrho_\varepsilon(y)\,d\sigma(y),
\]
because $\int\varrho_\varepsilon=1$ and $\mathbf1_{A''}H=H$ on
$\{\zeta-y:|y|\le\varepsilon\}\subseteq A''$; so
$|H_\varepsilon(\zeta)-H(\zeta)|\le\sup\{|H(\zeta')-H(\zeta'')|:
\zeta',\zeta''\in A'',\ |\zeta'-\zeta''|\le\varepsilon\}$, which tends to $0$
with $\varepsilon$ by uniform continuity.  Since
$\{re^{i\theta}:r\in[c,d]\}\subseteq A^\flat$ and $|e^{i(n-k)\theta}|=1$, it
follows that $g_\varepsilon\to a_k$ uniformly on $[c,d]$; in particular
$g_\varepsilon(c)\to a_k(c)$ and $g_\varepsilon(d)\to a_k(d)$.

Next, $\|W_{j,\varepsilon}-\mathbf1_{A''}D_j\|_{L^{1}(\mathbb{C})}\to0$, since
$g*\varrho_\varepsilon\to g$ in $L^{1}(\mathbb{C})$ for every
$g\in L^{1}(\mathbb{C})$ by \noderef{MollificationL1Convergence}.  On $A^\flat$ the coefficients
$\zeta_1/|\zeta|,\zeta_2/|\zeta|$ are bounded by $1$ and
$\mathbf1_{A''}D_j=D_j$, so
\[
  \int_{A^\flat}|R_\varepsilon-R|\,d\sigma
  \le\sum_{j=0}^1\|W_{j,\varepsilon}-\mathbf1_{A''}D_j\|_{L^{1}(\mathbb{C})}
  \longrightarrow0 .
\]
Passing to polar coordinates exactly as in (2), with $s\ge c-\delta>0$,
\[
  \Bigl|\int_c^d\Bigl(\frac1{2\pi}\int_0^{2\pi}
    (R_\varepsilon-R)(se^{i\theta})e^{i(n-k)\theta}d\theta\Bigr)ds\Bigr|
  \le\frac1{2\pi(c-\delta)}\int_{A^\flat}|R_\varepsilon-R|\,d\sigma
  \longrightarrow0 .
\]
So the right-hand side of (6) converges to
$\int_c^d\widetilde A_k(s)\,ds$, which by (3) equals $\int_c^dA_k(s)\,ds$.
Letting $\varepsilon\to0$ in (6) therefore gives
$a_k(d)-a_k(c)=\int_c^dA_k(s)\,ds$ whenever $c\le d$.  For $d<c$ apply this to
the pair $(d,c)$ and use $\int_c^d=-\int_d^c$; for $c=d$ both sides vanish.
This is (b).

\emph{Step 7: proof of (c).}  Fix $c_0\in(a,b)$.  By (a) the function $A_k$ is
integrable on every closed subinterval of $(a,b)$, hence locally integrable
there, and by (b) $a_k(r)=a_k(c_0)+\int_{c_0}^rA_k(s)\,ds$ for every
$r\in(a,b)$.  The indefinite integral of a locally integrable function on an
interval is differentiable at almost every point of that interval, with
derivative the integrand (the Lebesgue differentiation theorem in one
variable).  Hence for almost every $r\in(a,b)$ the function $a_k$ is
differentiable at $r$ with $a_k'(r)=A_k(r)$, which is (c).

\emph{Step 8: proof of (d).}  Fix $c,d\in(a,b)$ with $c\le d$ and keep the
notation of Steps 3--6.  For $r\in[c,d]$ the function
$\theta\mapsto H_\varepsilon(re^{i\theta})$ is $C^\infty$ on $\mathbb{R}$, and
the curve $\theta\mapsto re^{i\theta}$ has velocity $ire^{i\theta}=i\xi$ at
$\xi=re^{i\theta}$, so by the chain rule and \noderef{HAng}
\[
  \frac{\partial}{\partial\theta}\bigl[H_\varepsilon(re^{i\theta})\bigr]
  =\mathrm{d}H_\varepsilon(re^{i\theta})[i\,re^{i\theta}]
  =(H_\varepsilon)_\theta(re^{i\theta}).
\]
Both $\theta\mapsto H_\varepsilon(re^{i\theta})$ and
$\theta\mapsto e^{i(n-k)\theta}$ are $C^1$ on $[0,2\pi]$ and take the same
values at $0$ and at $2\pi$, the first because $e^{i\theta}$ is $2\pi$-periodic
and the second because $n-k$ is an integer.  So integration by parts on
$[0,2\pi]$ has vanishing boundary term and gives
\[
  \frac1{2\pi}\int_0^{2\pi}(H_\varepsilon)_\theta(re^{i\theta})
    e^{i(n-k)\theta}\,d\theta
  =-i(n-k)\frac1{2\pi}\int_0^{2\pi}H_\varepsilon(re^{i\theta})
    e^{i(n-k)\theta}\,d\theta
  =i(k-n)\,g_\varepsilon(r).
\]
By (5) the left-hand side is
$\frac1{2\pi}\int_0^{2\pi}S_\varepsilon(re^{i\theta})e^{i(n-k)\theta}d\theta$,
so
\begin{equation}\tag{7}
  \frac1{2\pi}\int_0^{2\pi}S_\varepsilon(re^{i\theta})
    e^{i(n-k)\theta}\,d\theta=i(k-n)g_\varepsilon(r)
  \qquad(r\in[c,d],\ 0<\varepsilon\le\delta).
\end{equation}
Now let $\varepsilon\to0$ in (7), comparing the two sides in $L^{1}[c,d]$.  On
$A^\flat$ the coefficients of $S$ and $S_\varepsilon$ are bounded by
$d+\delta$, so as in Step 6
\[
  \int_c^d\Bigl|\frac1{2\pi}\int_0^{2\pi}(S_\varepsilon-S)(se^{i\theta})
    e^{i(n-k)\theta}d\theta\Bigr|ds
  \le\frac{d+\delta}{2\pi(c-\delta)}
    \sum_{j=0}^1\|W_{j,\varepsilon}-\mathbf1_{A''}D_j\|_{L^{1}(\mathbb{C})}
  \longrightarrow0,
\]
so the left-hand side of (7) converges in $L^{1}[c,d]$ to $\widetilde B_k$,
which equals $B_k$ almost everywhere on $[c,d]$ by (3).  The right-hand side
converges to $i(k-n)a_k$ uniformly on $[c,d]$ by Step 6, hence in
$L^{1}[c,d]$.  Two $L^{1}[c,d]$ limits of one family agree almost everywhere,
so $B_k(r)=i(k-n)a_k(r)$ for almost every $r\in[c,d]$.

Finally, $(a,b)=\bigcup_{m\ge1}[c_m,d_m]$ with
$c_m=a+\tfrac{b-a}{m+2}$ and $d_m=b-\tfrac{b-a}{m+2}$, each contained in
$(a,b)$; a countable union of null sets is null, so the exceptional sets
assemble and (d) holds for almost every $r\in(a,b)$.

\emph{Step 9: proof of (e).}  Step 2, applied on the interval $[c_m,d_m]$,
established that for almost every $s\in[c_m,d_m]$ the integrand of
$\widetilde A_k(s)$ is integrable in $\theta$ and that the integrand of
$A_k(s)$ is integrable with it, and likewise for $\widetilde B_k$ and $B_k$;
those integrands are $\theta\mapsto R(se^{i\theta})e^{i(n-k)\theta}$ and
$\theta\mapsto H_r(se^{i\theta})e^{i(n-k)\theta}$, and the two angular ones.
Since $|e^{i(n-k)\theta}|=1$ and $\theta\mapsto e^{i(n-k)\theta}$ is
continuous, a function is integrable on $[0,2\pi]$ if and only if its product
with that factor is.  So for almost every $s\in[c_m,d_m]$ the functions
$\theta\mapsto H_r(se^{i\theta})$ and $\theta\mapsto H_\theta(se^{i\theta})$
are integrable on $[0,2\pi]$, and assembling the countably many exceptional
sets over the exhaustion above gives (e).
\end{proof}
